The single-machine comparison \acrshort{rq}1 poses resolves to a division of outcomes rather than a single superior stack. On wall-clock query time a warm, indexed PostGIS held the lowest run-time in every single-machine cell, and the GeoParquet and DuckDB path did not turn its decoupled read model into a run-time win at the sub-second scale these queries occupy. Network transfer was a per-architecture signature rather than a single ranking, because the bytes crossing the host boundary denote a scanned input for DuckDB and a returned result set for PostGIS, and the larger mover reversed by query pattern. At the single-node scale tested, that transfer did not reach operational cost, because the egress term was zero for every single-node configuration. DuckDB nonetheless recorded the lowest cost on two of the small-tier query patterns, despite trailing on run-time there.

For the national-scale spatial join, the answer to \acrshort{rq}2 turns on the size of the work. At the large tier neither single-node engine completed the join while the distributed broadcast join on Apache Sedona did, so distribution at that scale is a matter of feasibility rather than speed. Where both paths completed, the broadcast join scaled with added workers, yet its speedup stayed below the ideal-linear reference and fell further short as workers grew, the trace of a fixed coordination overhead and a partition load that more workers cannot drain. For a small join the comparison reversed: across the 2--16 workers tested the broadcast join approached the PostGIS minimum without overtaking it, so a tuned single node remained the better choice. From the medium tier, by contrast, the broadcast join led the single-node engines at every worker count. The choice of join strategy was the largest single lever, with the partitioned join running roughly two orders of magnitude slower than the broadcast join where it completed and exhausting executor memory at the medium and large tiers.

Read together, as \acrshort{rq}3 asks, those orderings did not hold a consistent winner. The leading configuration depended jointly on the query pattern, the dataset tier, and the outcome dimension, and changing any one of them could change it, so a single best engine is an ill-posed question. The practitioner guidance that follows is therefore conditional. For single-machine analytical queries, a tuned single-node PostGIS is the configuration to reach for, a lead that tracks the native spatial index it exposes rather than the engine in the abstract. Distribution earns its added operational complexity once the working set outgrows a single machine. DuckDB's value is operational, resting on its low cost and predictable read footprint rather than on latency.

The work contributed the first reported comparison of a distributed engine, Apache Sedona, against single-node PostGIS and DuckDB on a national-scale spatial join. It also reported network transfer jointly with run-time and operational cost, a dimension absent from almost all spatial benchmarks. Methodologically, it reported effect sizes and bootstrapped confidence intervals and adopted a rank-based omnibus test with post-hoc correction, practices rarely found in geospatial benchmarking. These findings are bounded. They rest on a single cloud provider, on vector building polygons rather than raster or other feature types, on read performance alone, on a single client issuing queries in sequence, and on a distributed sweep of 2--16 workers. The orderings that follow from format and engine structure are expected to reach beyond those conditions, while the absolute times and costs, priced from provider- and region-specific rates, are not. One bound shapes the work that follows: the single-node run-time comparison varied indexing capability and storage locality together with the engine, and the GeoParquet path was left unindexed beyond row-group pruning.
